\chapter{Insiemi, Funzioni e Grafici}
\label{chap:insiemi_funzioni}

Lo studio dell'Analisi Matematica poggia su fondamenta concettuali precise: la struttura dei numeri reali, la formalizzazione del concetto di funzione e la traduzione geometrica delle relazioni analitiche attraverso il piano cartesiano. Prima di affrontare il calcolo infinitesimale (limiti, derivate ed integrali), è indispensabile acquisire una piena padronanza del linguaggio insiemistico, delle trasformazioni funzionali e dei sistemi di coordinate.

In questo capitolo sviluppiamo una trattazione organica e lineare che parte dalla costruzione logica dei sistemi numerici, analizza il comportamento qualitativo e le simmetrie delle funzioni reali, approfondisce i sistemi di coordinate (cartesiane e polari), risolve i paradossi classici relativi a radicali e potenze, e introduce la trigonometria iperbolica con i suoi profondi legami geometrici.

% ==============================================================================
\section{Insiemi Numerici e Struttura dei Numeri Reali}
\label{sec:insiemi_numerici}
% ==============================================================================

Il percorso logico che conduce alla costruzione dell'insieme dei numeri reali $\R$ risponde all'esigenza storica e formale di dotare la retta geometrica di una struttura algebrica priva di ``lacune'' o buchi.

\subsection{La Catena delle Inclusioni Numeriche}
Partendo dal conteggio degli elementi discreti, si definisce la successione di ampliamenti algebrici:
\[
    \N \subset \Z \subset \Q \subset \R \subset \C
\]

\begin{definizione}{Insiemi Numerici Fondamentali}{insiemi_numerici}
\begin{enumerate}[label=\textbf{\arabic*.}]
    \item \textbf{Numeri Naturali} ($\N$): $\N = \{0, 1, 2, 3, \dots\}$. È la struttura minima dotata dell'operazione di successore e dell'assioma di induzione matematica.
    \item \textbf{Numeri Interi Relativi} ($\Z$): $\Z = \{0, \pm 1, \pm 2, \pm 3, \dots\}$. Costituisce un anello commutativo unitario in cui ogni elemento ammette opposto rispetto all'addizione.
    \item \textbf{Numeri Razionali} ($\Q$): $\Q = \left\{ \frac{p}{q} \;\middle|\; p \in \Z, \, q \in \N_{>0}, \, \gcd(p, q) = 1 \right\}$. Rappresenta un campo ordinato ed archimedeo nel quale ogni elemento non nullo possiede reciproco rispetto alla moltiplicazione.
    \item \textbf{Numeri Reali} ($\R$): Campo ordinato, archimedeo e \textbf{completo}, ottenuto come completamento topologico di $\Q$.
    \item \textbf{Numeri Complessi} ($\C$): Chiusura algebrica di $\R$, definita come $\C = \{x + iy \mid x, y \in \R, i^2 = -1\}$.
\end{enumerate}
\end{definizione}

Sebbene l'insieme dei numeri razionali $\Q$ sia denso in sé (ossia tra due razionali distinti esiste sempre un terzo numero razionale, per esempio la media aritmetica $\frac{q_1+q_2}{2}$), esso presenta una densità discontinua. Storicamente, la scoperta dei pitagorici dell'incommensurabilità della diagonale del quadrato unitario dimostrò che non esiste alcun numero $x \in \Q$ tale che $x^2 = 2$. Da ciò scaturisce la necessità dell'assioma di completezza.

\subsection{L'Assioma di Completezza di Dedekind}
La retta reale $\R$ si distingue da $\Q$ per una proprietà topologica fondamentale: ogni sottoinsieme non vuoto e superiormente limitato ammette un estremo superiore appartenente a $\R$.

\begin{definizione}{Estremo Superiore ed Inferiore}{sup_inf}
Sia $E \subseteq \R$ un sottoinsieme non vuoto.
\begin{itemize}
    \item Un elemento $M \in \R$ è un \textbf{maggiorante} di $E$ se $\forall x \in E, x \le M$. Se l'insieme dei maggioranti non è vuoto, $E$ si dice superiormente limitato.
    \item Si definisce \textbf{estremo superiore} di $E$, denotato con $\sup E$, il minimo dei maggioranti di $E$:
    \[
        L = \sup E \iff \begin{cases}
            \forall x \in E, \quad x \le L & \text{(L è un maggiorante)} \\
            \forall \eps > 0, \; \exists x_\eps \in E : x_\eps > L - \eps & \text{(L è il minimo maggiorante)}
        \end{cases}
    \]
    \item Analogamente, si definisce \textbf{estremo inferiore} $\inf E$ come il massimo dei minoranti.
\end{itemize}
Se l'estremo superiore appartiene all'insieme $E$, esso prende il nome di \textbf{massimo assoluto} ($\max E = \sup E$).
\end{definizione}

\begin{teorema}{Assioma di Completezza (di Dedekind)}{completezza_dedekind}
Ogni sottoinsieme non vuoto $E \subset \R$ superiormente limitato ammette estremo superiore finito in $\R$: $\sup E \in \R$.
\end{teorema}

\begin{osservazione}[Conseguenze della Completezza]
L'assioma di completezza garantisce la validità di tutti i grandi teoremi dell'Analisi 1: il Teorema di Weierstrass sui compatti, il Teorema di esistenza degli zeri per funzioni continue, la convergenza delle successioni monotone limitate e la convergenza di tutte le successioni di Cauchy.
\end{osservazione}

% ==============================================================================
\section{Il Concetto Moderno di Funzione e Proprietà Strutturali}
\label{sec:funzioni_proprieta}
% ==============================================================================

Una funzione matematica non è una semplice formula algebrica, ma una terna ordinata formata da un insieme di partenza (dominio), un insieme di arrivo (codominio) e una legge univoca di associazione.

\begin{definizione}{Funzione Reale di Variabile Reale}{funzione_reale}
Dati due sottoinsiemi $A, B \subseteq \R$, una \textbf{funzione} $f \colon A \to B$ è una corrispondenza che associa a ciascun elemento $x \in A$ uno ed un solo elemento $y = f(x) \in B$.
\begin{itemize}
    \item $A = \dom(f)$ è l'\textbf{insieme di definizione} o \textbf{dominio naturale}.
    \item $B = \codom(f)$ è il \textbf{codominio}.
    \item $\Imm(f) = f(A) = \{y \in B \mid \exists x \in A : f(x) = y\}$ è l'\textbf{immagine} di $f$.
    \item Il \textbf{grafico} di $f$ è il sottoinsieme del piano cartesiano $\R^2$:
    \[
        \operatorname{Graph}(f) = \big\{ (x, y) \in \R^2 \;\big|\; x \in A, \, y = f(x) \big\}
    \]
\end{itemize}
\end{definizione}

\subsection{Iniettività, Suriettività e Funzione Inversa}
Le proprietà qualitative di una funzione determinano la possibilità di invertire la relazione di corrispondenza.

\begin{definizione}{Classificazione delle Applicazioni}{classificazione_applicazioni}
Sia $f \colon A \to B$:
\begin{enumerate}
    \item \textbf{Iniettiva}: elementi distinti del dominio hanno immagini distinte:
    \[
        \forall x_1, x_2 \in A, \quad x_1 \neq x_2 \implies f(x_1) \neq f(x_2) \quad \left(\text{o eq. } f(x_1) = f(x_2) \implies x_1 = x_2\right)
    \]
    \emph{Interpretazione geometrica}: ogni retta orizzontale $y = y_0$ interseca il grafico di $f$ in \textbf{al più un punto}.
    \item \textbf{Suriettiva}: ogni elemento del codominio è immagine di almeno un elemento del dominio:
    \[
        \forall y \in B, \quad \exists x \in A : f(x) = y \quad \iff \quad \Imm(f) = B
    \]
    \emph{Interpretazione geometrica}: ogni retta orizzontale $y = y_0$ con $y_0 \in B$ interseca il grafico in \textbf{almeno un punto}.
    \item \textbf{Biettiva} (o biunivoca): $f$ è contemporaneamente iniettiva e suriettiva. Ogni retta orizzontale $y = y_0$ interseca il grafico in \textbf{uno ed un solo punto}.
\end{enumerate}
\end{definizione}

\begin{teorema}{Esistenza e Unicità della Funzione Inversa}{funzione_inversa}
Sia $f \colon A \to B$ una funzione biettiva. Allora esiste un'unica funzione, detta \textbf{funzione inversa} e denotata con $f^{-1} \colon B \to A$, tale che:
\[
    f^{-1}(f(x)) = x \quad \forall x \in A, \qquad f(f^{-1}(y)) = y \quad \forall y \in B
\]
\end{teorema}

\begin{dimostrazione}
Poiché $f$ è suriettiva su $B$, per ogni $y \in B$ l'insieme della controimmagine $f^{-1}(\{y\}) = \{x \in A \mid f(x) = y\}$ è non vuoto. Poiché $f$ è iniettiva, tale insieme contiene esattamente un unico elemento $x \in A$. Definendo $g(y) = x$, la corrispondenza $g \colon B \to A$ è univocamente determinata ed è la funzione inversa cercata.
\end{dimostrazione}

\begin{osservazione}[Simmetria Geometrica del Grafico Inverso]
Il punto $(x_0, y_0) \in \operatorname{Graph}(f)$ se e solo se $y_0 = f(x_0)$, il che equivale a $x_0 = f^{-1}(y_0)$, ossia $(y_0, x_0) \in \operatorname{Graph}(f^{-1})$.
Geometricamente, il punto $(y_0, x_0)$ si ottiene da $(x_0, y_0)$ mediante una \textbf{riflessione assiale rispetto alla bisettrice del I e III quadrante} ($y = x$).
\end{osservazione}

% ==============================================================================
\section{Algebra delle Trasformazioni Geometriche dei Grafici}
\label{sec:trasformazioni_grafici_dettagli}
% ==============================================================================

Conoscere il grafico qualitativo delle funzioni elementari ($x^n, e^x, \ln x, \sen x, \cos x, \dots$) consente di tracciare con rapidità e precisione grafici complessi applicando una sequenza logica di trasformazioni geometriche elementari.

\begin{metodo}[Catalogo delle Trasformazioni Geometriche]
Sia noto il grafico della funzione $y = f(x)$:
\begin{enumerate}
    \item \textbf{Traslazione verticale}: $y = f(x) + c$. Se $c > 0$ il grafico si sposta verso l'alto di $c$ unità; se $c < 0$ si sposta verso il basso di $|c|$ unità.
    \item \textbf{Traslazione orizzontale}: $y = f(x - c)$. Se $c > 0$ il grafico trasla verso destra di $c$; se $c < 0$ verso sinistra di $|c|$.
    \item \textbf{Riflessione rispetto all'asse $x$}: $y = -f(x)$. Le quote positive diventano negative e viceversa.
    \item \textbf{Riflessione rispetto all'asse $y$}: $y = f(-x)$. Il semipiano destro viene ribaltato a sinistra e viceversa. Se $f(-x) = f(x)$ la funzione è \textbf{pari} (simmetria assiale rispetto a $x = 0$); se $f(-x) = -f(x)$ la funzione è \textbf{dispari} (simmetria centrale rispetto all'origine $O$).
    \item \textbf{Dilatazione / Contrazione}:
    \begin{itemize}
        \item $y = a f(x)$ ($a > 0$): riscalamento verticale di fattore $a$.
        \item $y = f(k x)$ ($k > 0$): compressione orizzontale se $k > 1$, dilatazione orizzontale se $0 < k < 1$.
    \end{itemize}
    \item \textbf{Valore Assoluto Esterno}: $y = |f(x)|$. Le porzioni di grafico con $f(x) \ge 0$ rimangono inalterate, mentre le porzioni in cui $f(x) < 0$ vengono ribaltate simmetricamente verso l'alto nel semipiano $y \ge 0$.
    \item \textbf{Valore Assoluto Interno}: $y = f(|x|)$. Si conserva il ramo per $x \ge 0$ e lo si riflette a specchio a sinistra per $x < 0$, ottenendo una funzione automaticamente pari.
\end{enumerate}
\end{metodo}

\begin{esame}[Attenzione all'Ordine delle Trasformazioni Composte]
L'ordine di applicazione delle trasformazioni è vincolante e non commutativo.
Ad esempio, per tracciare $y = f(2x + 4)$:
\begin{itemize}
    \item \emph{Procedura corretta}: riscrivere come $y = f(2(x + 2))$. Prima si applica la compressione orizzontale di fattore 2 ($x \mapsto 2x$), e successivamente si trasla a sinistra di 2 unità ($x \mapsto x + 2$).
    \item Se si traslasse prima di 4 e poi si comprimesse, si otterrebbe $f(2x + 4)$ solo applicando la compressione all'intera variabile traslata.
\end{itemize}
\end{esame}

% ==============================================================================
\section{Sistemi di Coordinate nel Piano: Cartesiane e Polari}
\label{sec:coordinate_polari_teoria}
% ==============================================================================

Nel piano euclideo, la posizione di un punto può essere descritta specificando le distanze ortogonali rispetto a due assi di riferimento (coordinate cartesiane) oppure la distanza dal polo e l'angolo di orientamento rispetto a un asse polare (coordinate polari).

\begin{definizione}{Coordinate Polari nel Piano}{coordinate_polari_def}
Fissato un polo $O \equiv (0, 0)$ e un semiasse polare coincidente con il semiasse positivo delle ascisse $x \ge 0$:
\begin{itemize}
    \item $\rho \in [0, +\infty)$ è il \textbf{raggio polare} (distanza euclidea dal polo $O$).
    \item $\theta \in (-\pi, \pi]$ (convenzione dell'argomento principale) oppure $\theta \in [0, 2\pi)$ è l'\textbf{anomalia} o angolo polare orientato.
\end{itemize}
\end{definizione}

\begin{figure}[htbp]
    \centering
    \begin{tikzpicture}[scale=1.2, >=Stealth]
    % Griglia leggera di sfondo
    \draw[step=0.5cm, gray!20, very thin] (-3.2,-2.2) grid (3.5,3.2);
    
    % Assi cartesiani
    \draw[->, thick, navyblue] (-3.2,0) -- (3.5,0) node[right] {$x$};
    \draw[->, thick, navyblue] (0,-2.2) -- (0,3.2) node[above] {$y$};
    \node[below left, font=\small] at (0,0) {$O$};

    % Coordinate del punto P
    \def\r{2.8}
    \def\ang{40}
    \coordinate (O) at (0,0);
    \coordinate (P) at ({\r*cos(\ang)}, {\r*sin(\ang)});
    \coordinate (Px) at ({\r*cos(\ang)}, 0);
    \coordinate (Py) at (0, {\r*sin(\ang)});

    % Triangolo rettangolo riempito
    \fill[coolblue!10] (O) -- (Px) -- (P) -- cycle;
    
    % Proiezioni ortogonali tratteggiate
    \draw[dashed, coolblue, thick] (P) -- (Px) node[below, black] {$x = \rho \cos\theta$};
    \draw[dashed, coolblue, thick] (P) -- (Py) node[left, black] {$y = \rho \sin\theta$};

    % Angolo retto
    \draw[gray!70, thin] ({\r*cos(\ang)-0.25}, 0) |- ({\r*cos(\ang)}, 0.25);
    \fill[gray!70] ({\r*cos(\ang)-0.12}, 0.12) circle (0.8pt);

    % Vettore raggio rho
    \draw[->, very thick, crimsonred] (O) -- (P) node[midway, above left, font=\bfseries] {$\rho = \sqrt{x^2+y^2}$};

    % Arco angolo theta
    \draw[->, thick, forestgreen] (0.8,0) arc (0:\ang:0.8);
    \node[forestgreen, font=\bfseries] at (20:1.1) {$\theta$};

    % Punto P
    \fill[navyblue] (P) circle (2.2pt);
    \node[above right, font=\small] at (P) {$P(x,y) \equiv P(\rho, \theta)$};

    % Etichette quadranti e formule di correzione
    \node[font=\footnotesize, align=center, fill=white, inner sep=1pt, opacity=0.9, text opacity=1] at (2.2, 2.2) {\textbf{I Quadrante}\\$\theta = \arctan(y/x)$};
    \node[font=\footnotesize, align=center, fill=white, inner sep=1pt, opacity=0.9, text opacity=1] at (-1.8, 2.2) {\textbf{II Quadrante}\\$\theta = \arctan(y/x) + \pi$};
    \node[font=\footnotesize, align=center, fill=white, inner sep=1pt, opacity=0.9, text opacity=1] at (-1.8, -1.5) {\textbf{III Quadrante}\\$\theta = \arctan(y/x) - \pi$};
    \node[font=\footnotesize, align=center, fill=white, inner sep=1pt, opacity=0.9, text opacity=1] at (2.2, -1.5) {\textbf{IV Quadrante}\\$\theta = \arctan(y/x)$};
\end{tikzpicture}

    \caption{Rappresentazione geometrica della conversione tra coordinate cartesiane $(x,y)$ e coordinate polari $(\rho, \theta)$, con evidenziazione del triangolo rettangolo fondamentale e tabella analitica delle correzioni per quadrante.}
    \label{fig:coord_polari_schema}
\end{figure}

\subsection{Dimostrazione Rigorosa delle Formule di Conversione}
\paragraph{1. Da Polari a Cartesiane:}
Considerando il triangolo rettangolo di ipotenusa $\rho$ avente per cateti le proiezioni ortogonali lungo gli assi, le definizioni trigonometriche elementari porgono direttamente:
\[
    x = \rho \cos \theta, \qquad y = \rho \sen \theta
\]

\paragraph{2. Da Cartesiane a Polari:}
Sommando i quadrati delle relazioni cartesiane:
\[
    x^2 + y^2 = \rho^2 \cos^2 \theta + \rho^2 \sen^2 \theta = \rho^2 (\cos^2 \theta + \sen^2 \theta) = \rho^2 \implies \rho = \sqrt{x^2 + y^2} \ge 0
\]
Per quanto riguarda l'angolo $\theta$, per ogni punto con $x \neq 0$ si ha $\frac{y}{x} = \frac{\rho \sen \theta}{\rho \cos \theta} = \tg \theta$.
Tuttavia, la funzione arcotangente $\arctg \colon \R \to (-\pi/2, \pi/2)$ restituisce angoli strettamente confinati nel I e IV quadrante. Qualora il punto $(x, y)$ appartenga al II o al III quadrante ($x < 0$), l'angolo effettivo differisce di $\pi$ dall'uscita dell'arcotangente:
\[
    \theta = \begin{cases}
        \arctg\left(\frac{y}{x}\right) & \text{se } x > 0, \; y \in \R \quad (\text{I e IV quadrante}) \\[1ex]
        \arctg\left(\frac{y}{x}\right) + \pi & \text{se } x < 0, \; y \ge 0 \quad (\text{II quadrante}) \\[1ex]
        \arctg\left(\frac{y}{x}\right) - \pi & \text{se } x < 0, \; y < 0 \quad (\text{III quadrante, conv. } (-\pi, \pi]) \\[1ex]
        +\frac{\pi}{2} & \text{se } x = 0, \; y > 0 \\[1ex]
        -\frac{\pi}{2} & \text{se } x = 0, \; y < 0
    \end{cases}
\]

% ==============================================================================
\section{Teoria Assiomatica dei Radicali e delle Potenze}
\label{sec:potenze_radicali_teoria}
% ==============================================================================

Una delle fonti più frequenti di errore nel calcolo differenziale risiede nell'errata manipolazione dei radicali con indice pari e nella falsa estensione delle potenze ad esponente frazionario o reale su basi negative.

\subsection{Radici Aritmetiche vs Algebriche e l'Identità $\sqrt{x^2} = |x|$}
In Analisi Matematica, il simbolo $\sqrt{a}$ indica la **radice aritmetica positiva**.

\begin{definizione}{Radice Aritmetica $n$-esima}{radice_aritmetica_def}
Dato $n \in \N_{\ge 1}$ e $a \in \R_{\ge 0}$, si definisce \textbf{radice aritmetica $n$-esima} di $a$, indicata con $\sqrt[n]{a}$, l'unico numero reale $b \ge 0$ tale che $b^n = a$.
\end{definizione}

\begin{teorema}{Identità Fondamentale del Valore Assoluto}{radice_quadrata_modulo}
Per ogni numero reale $x \in \R$:
\[
    \sqrt{x^2} = |x|
\]
\end{teorema}

\begin{dimostrazione}
Per ogni $x \in \R$, $x^2 \ge 0$, dunque $\sqrt{x^2}$ è ben definita. Per definizione di radice aritmetica, $\sqrt{x^2}$ è l'unico valore $b \ge 0$ tale che $b^2 = x^2$.
Verifichiamo che $b = |x|$ soddisfa entrambi i requisiti:
1. $|x| \ge 0$ per definizione di valore assoluto;
2. $(|x|)^2 = (\pm x)^2 = x^2$.
Per l'unicità della radice aritmetica, ne consegue che $\sqrt{x^2} = |x|$.
\end{dimostrazione}

\begin{esame}[Trabocchetto d'Esame: Estrazione nei Limiti a $-\infty$]
Uno degli errori più penalizzati consiste nello scrivere $\sqrt{x^2} = x$ quando $x \to -\infty$.
Infatti, poiché per $x < 0$ si ha $|x| = -x$:
\[
    \lim_{x \to -\infty} \frac{\sqrt{4x^2 + 1}}{x} = \lim_{x \to -\infty} \frac{\sqrt{x^2(4 + x^{-2})}}{x} = \lim_{x \to -\infty} \frac{|x|\sqrt{4 + x^{-2}}}{x} = \lim_{x \to -\infty} \frac{-x\sqrt{4 + x^{-2}}}{x} = -2 \quad (\neq +2)
\]
\end{esame}

\subsection{Gerarchia delle Potenze e il Paradosso delle Basi Negative}
L'elevamento a potenza si costruisce rigorosamente in quattro passaggi successivi:
\begin{enumerate}
    \item \textbf{Esponente naturale} ($n \in \N_{>0}$): $a^n = \underbrace{a \cdot a \cdots a}_{n \text{ volte}}$ per ogni $a \in \R$.
    \item \textbf{Esponente intero} ($m \in \Z$): $a^0 = 1$ ($a \neq 0$), $a^{-n} = \frac{1}{a^n}$ ($a \neq 0, n \in \N_{>0}$).
    \item \textbf{Esponente razionale} ($q = \frac{m}{n} \in \Q$ con $n > 0$): $a^{m/n} = \sqrt[n]{a^m}$, definita per **$a > 0$** (e per $a = 0$ se $m/n > 0$).
    \item \textbf{Esponente reale irrazionale} ($x \in \R \setminus \Q$): definito per continuità mediante l'estremo superiore delle potenze razionali approssimanti, oppure tramite l'identità analitica fondamentale:
    \[
        a^x := e^{x \ln a} \qquad (a > 0)
    \]
\end{enumerate}

\paragraph{Perché le potenze razionali non possono essere definite per basi negative?}
Se ammettessimo basi negative, nascerebbe una contraddizione insanabile legata all'equivalenza delle frazioni razionali.
Consideriamo $a = -8$ e l'esponente $1/3 = 2/6 \in \Q$:
\begin{align*}
    (-8)^{1/3} &= \sqrt[3]{-8} = -2 \\
    (-8)^{2/6} &= \sqrt[6]{(-8)^2} = \sqrt[6]{64} = +2
\end{align*}
Poiché $-2 \neq +2$, l'operazione non sarebbe ben definita sulla classe di equivalenza dei numeri razionali. Inoltre verrebbero meno le proprietà algebriche fondamentali $(a^p)^q = a^{pq}$:
\[
    \left((-1)^2\right)^{1/2} = 1^{1/2} = 1 \neq (-1)^{2 \cdot (1/2)} = (-1)^1 = -1
\]
Pertanto, per garantire la coerenza logica e la continuità, la funzione esponenziale $x \mapsto a^x$ è rigorosamente definita solo per $a > 0$.

% ==============================================================================
\section{Trigonometria Iperbolica e Geometria dell'Iperbole Equilatera}
\label{sec:trigonometria_iperbolica_completa}
% ==============================================================================

In perfetta analogia con le funzioni trigonometriche circolari (generate dalla circonferenza $X^2 + Y^2 = 1$), le **funzioni iperboliche** scaturiscono dalla parametrizzazione dell'iperbole equilatera $X^2 - Y^2 = 1$.

\begin{definizione}{Funzioni Iperboliche}{funzioni_iperboliche_def}
Per ogni $x \in \R$, si definiscono:
\[
    \cosh x = \frac{e^x + e^{-x}}{2}, \qquad \sinh x = \frac{e^x - e^{-x}}{2}, \qquad \tanh x = \frac{\sinh x}{\cosh x} = \frac{e^x - e^{-x}}{e^x + e^{-x}}
\]
\end{definizione}

\begin{figure}[htbp]
    \centering
    \begin{tikzpicture}
\begin{axis}[
    axis lines=middle,
    grid=both,
    grid style={line width=.1pt, draw=gray!20},
    major grid style={line width=.2pt,draw=gray!40},
    xlabel={$x$},
    ylabel={$y$},
    xmin=-3.2, xmax=3.2,
    ymin=-3.5, ymax=4.2,
    domain=-3:3,
    samples=150,
    axis line style={thick, navyblue},
    legend style={at={(0.03,0.97)}, anchor=north west, font=\small, fill=white, fill opacity=0.9, draw=gray!50},
    legend cell align={left},
    width=11cm,
    height=8.5cm
]
    % Asintoti orizzontali per tanh(x)
    \addplot[dashed, crimsonred!80, thick, domain=-3.2:3.2] {1};
    \addplot[dashed, crimsonred!80, thick, domain=-3.2:3.2] {-1};
    \node[crimsonred, font=\footnotesize, anchor=south east] at (axis cs:3.1,1) {$y = 1$};
    \node[crimsonred, font=\footnotesize, anchor=north east] at (axis cs:3.1,-1) {$y = -1$};

    % Funzione esponenziale e^x / 2 asintoto per cosh e sinh
    \addplot[loosely dotted, gray!80, thick, domain=-0.5:2.2] {0.5*exp(x)};
    \node[gray!80, font=\footnotesize, anchor=west] at (axis cs:2.1, 4.0) {$y = \frac{1}{2}e^x$};

    % Cosh(x)
    \addplot[very thick, coolblue] {cosh(x)};
    \addlegendentry{$\cosh(x) = \frac{e^x + e^{-x}}{2}$}

    % Sinh(x)
    \addplot[very thick, forestgreen] {sinh(x)};
    \addlegendentry{$\sinh(x) = \frac{e^x - e^{-x}}{2}$}

    % Tanh(x)
    \addplot[very thick, warmamber] {tanh(x)};
    \addlegendentry{$\tanh(x) = \frac{\sinh(x)}{\cosh(x)}$}

    % Punti notevoli
    \fill[coolblue] (axis cs:0,1) circle (2pt) node[above left, font=\footnotesize] {$(0,1)$};
    \fill[forestgreen] (axis cs:0,0) circle (2pt) node[below right, font=\footnotesize] {$(0,0)$};

\end{axis}
\end{tikzpicture}

    \caption{Grafici delle funzioni iperboliche $\cosh(x)$ (pari, curva catenaria), $\sinh(x)$ (dispari) e $\tanh(x)$ (dispari con asintoti orizzontali $y = \pm 1$), con il profilo asintotico esponenziale $y = \frac{1}{2}e^x$.}
    \label{fig:funzioni_iperboliche_plot}
\end{figure}

\begin{teorema}{Identità Fondamentale della Trigonometria Iperbolica}{identita_iperbolica}
Per ogni $x \in \R$:
\[
    \cosh^2 x - \sinh^2 x = 1
\]
\end{teorema}

\begin{dimostrazione}
Sviluppando direttamente i quadrati delle definizioni:
\begin{align*}
    \cosh^2 x &= \left(\frac{e^x + e^{-x}}{2}\right)^2 = \frac{e^{2x} + 2 + e^{-2x}}{4} \\
    \sinh^2 x &= \left(\frac{e^x - e^{-x}}{2}\right)^2 = \frac{e^{2x} - 2 + e^{-2x}}{4}
\end{align*}
Sottraendo membro a membro:
\[
    \cosh^2 x - \sinh^2 x = \frac{(e^{2x} + 2 + e^{-2x}) - (e^{2x} - 2 + e^{-2x})}{4} = \frac{4}{4} = 1
\]
\end{dimostrazione}

\begin{figure}[htbp]
    \centering
    \begin{tikzpicture}[scale=1.1, >=Stealth]
    % --- Riquadro Sinistro: Circonferenza Unitaria ---
    \begin{scope}[shift={(0,0)}]
        % Titolo
        \node[font=\bfseries\color{navyblue}] at (0, 2.6) {Circonferenza: $X^2 + Y^2 = 1$};
        
        % Assi
        \draw[->, thick, gray!80] (-2.2, 0) -- (2.2, 0) node[right] {$X$};
        \draw[->, thick, gray!80] (0, -2.2) -- (0, 2.2) node[above] {$Y$};
        
        % Settore circolare ombreggiato (Area = t/2)
        \def\angcircle{48}
        \fill[coolblue!25] (0,0) -- (1.8,0) arc (0:\angcircle:1.8) -- cycle;
        \node[coolblue!90!black, font=\bfseries] at (24:1.0) {Area $= \frac{t}{2}$};

        % Circonferenza
        \draw[very thick, coolblue] (0,0) circle (1.8);

        % Raggio e coordinate
        \coordinate (Oc) at (0,0);
        \coordinate (Pc) at ({\angcircle}:1.8);
        \draw[thick, crimsonred] (Oc) -- (Pc);
        \fill[crimsonred] (Pc) circle (2pt);
        \node[above right, crimsonred, font=\footnotesize] at (Pc) {$P(\cos t, \sin t)$};

        % Proiezioni
        \draw[dashed, gray] (Pc) -- ({\angcircle}:1.8 |- Oc) node[below, font=\scriptsize] {$\cos t$};
        \draw[dashed, gray] (Pc) -- (Pc -| Oc) node[left, font=\scriptsize] {$\sin t$};

        % Angolo t
        \draw[->, thick, forestgreen] (0.6,0) arc (0:\angcircle:0.6);
        \node[forestgreen, font=\footnotesize] at (24:0.8) {$t$};
    \end{scope}

    % --- Riquadro Destro: Iperbole Equilatera ---
    \begin{scope}[shift={(6.5,0)}]
        % Titolo
        \node[font=\bfseries\color{navyblue}] at (0, 2.6) {Iperbole Equilatera: $X^2 - Y^2 = 1$};
        
        % Assi
        \draw[->, thick, gray!80] (-1.2, 0) -- (3.2, 0) node[right] {$X$};
        \draw[->, thick, gray!80] (0, -2.2) -- (0, 2.2) node[above] {$Y$};

        % Asintoti Y = +- X
        \draw[dashed, gray!70] (-1.2, -1.2) -- (2.2, 2.2) node[above right, font=\tiny] {$Y=X$};
        \draw[dashed, gray!70] (-1.2, 1.2) -- (2.2, -2.2) node[below right, font=\tiny] {$Y=-X$};

        % Settore iperbolico ombreggiato (Area = t/2)
        % Usiamo coordinate parametriche per riempire: t_val = 1.1, scala = 1.2
        % x = a cosh(u), y = a sinh(u) con a = 1.4 (scala grafica)
        \def\a{1.3}
        \def\tmax{0.95}
        \fill[forestgreen!25] (0,0) -- ({\a*cosh(0)},0)
            \foreach \u in {0.05, 0.1, ..., \tmax} {
                -- ({\a*cosh(\u)}, {\a*sinh(\u)})
            }
            -- (0,0);
        \node[forestgreen!90!black, font=\bfseries] at (1.1, 0.4) {Area $= \frac{t}{2}$};

        % Ramo destro iperbole
        \draw[very thick, forestgreen, domain=-1.2:1.2, samples=100] 
            plot ({\a*cosh(\x)}, {\a*sinh(\x)});

        % Punto P e raggio
        \coordinate (Oh) at (0,0);
        \coordinate (Ph) at ({\a*cosh(\tmax)}, {\a*sinh(\tmax)});
        \draw[thick, crimsonred] (Oh) -- (Ph);
        \fill[crimsonred] (Ph) circle (2pt);
        \node[above right, crimsonred, font=\footnotesize] at (Ph) {$P(\cosh t, \sinh t)$};

        % Proiezioni
        \draw[dashed, gray] (Ph) -- ({\a*cosh(\tmax)}, 0) node[below, font=\scriptsize] {$\cosh t$};
        \draw[dashed, gray] (Ph) -- (0, {\a*sinh(\tmax)}) node[left, font=\scriptsize] {$\sinh t$};

        % Vertice
        \fill[black] (\a, 0) circle (1.5pt) node[below left, font=\scriptsize] {$(1,0)$};
    \end{scope}
\end{tikzpicture}

    \caption{Dualità geometrica: il cerchio unitario $X^2+Y^2=1$ parametrizzato da $(\cos t, \sen t)$ e l'iperbole equilatera $X^2-Y^2=1$ parametrizzata da $(\cosh t, \sinh t)$. In entrambi i casi, l'area del settore evidenziato vale esattamente $A = t/2$.}
    \label{fig:dualita_iperbole_cerchio}
\end{figure}

\subsection{Derivate e Significato Geometrico del Parametro $t$}
Le regole di derivazione delle funzioni iperboliche evidenziano la completa simmetria priva di segni alterni per il coseno:
\[
    \frac{\diff}{\dx} \sinh x = \cosh x, \qquad \frac{\diff}{\dx} \cosh x = +\sinh x, \qquad \frac{\diff}{\dx} \tanh x = \frac{1}{\cosh^2 x} = 1 - \tanh^2 x
\]

Inoltre, mentre nel cerchio goniometrico il parametro $t$ corrisponde sia alla lunghezza d'arco che al doppio dell'area del settore circolare ($A = t/2$), nell'iperbole equilatera $t$ rappresenta **esattamente il doppio dell'area del settore iperbolico** compreso tra l'origine, il vertice $(1, 0)$ e il punto $(\cosh t, \sinh t)$:
\[
    A(t) = \frac{1}{2}\cosh t \sinh t - \int_1^{\cosh t} \sqrt{u^2 - 1}\du = \frac{t}{2} \iff t = 2 A(t)
\]
Da questa profonda proprietà geometrica discende la denominazione delle funzioni inverse iperboliche:
\begin{itemize}
    \item Settore seno iperbolico: $\operatorname{settsh} x = \ln\left(x + \sqrt{x^2 + 1}\right), \quad x \in \R$
    \item Settore coseno iperbolico: $\operatorname{settch} x = \ln\left(x + \sqrt{x^2 - 1}\right), \quad x \ge 1$
    \item Settore tangente iperbolica: $\operatorname{setttgh} x = \frac{1}{2}\ln\left(\frac{1+x}{1-x}\right), \quad x \in (-1, 1)$
\end{itemize}
